\documentclass[a4paper,12pt]{article}
\usepackage{amsmath, amssymb}
\usepackage{geometry}
\usepackage{enumitem}
\usepackage{fancyhdr}
\usepackage{xcolor}
\usepackage{sectsty}
\usepackage{multicol}
\usepackage{graphicx}
\usepackage{float}
\def\inputGnumericTable{}
\usepackage[latin1]{inputenc}
\usepackage{fullpage}
\usepackage{color}
\usepackage{array}
\usepackage{longtable}
\usepackage{calc}
\usepackage{multirow}
\usepackage{hhline}
\usepackage{ifthen}
\geometry{top=1in, bottom=1.5cm, left=1.5cm, right=1.5cm}    % Header and Footer settings
\pagestyle{empty}
\sectionfont{\color{blue}}
\begin{document}

\thispagestyle{fancy}
\fancyhf{} % Clear default header and footer
\fancyhead[L]{% Left header
        
\includegraphics[width=8cm, height=1.7cm]{images/logo.png} % Adjust dimensions
}
\fancyhead[R]{% Right header
    Name: M.Thrinath Rao \\
    Batch: COMETFWC034 \\
    Date:  9 september 2025
}
\renewcommand{\headrulewidth}{0pt} % Remove header line
\fancyfoot[C]{\thepage} % Page number centered in footer
\vspace*{1cm}
\begin{center}
        {\LARGE \textbf{\textcolor{blue}{GATE 2010 EE, $52^{nd}$ Question Analysis}}}
\end{center}
\section*{\textbf{Question 52}}
\begin{center}
        \includegraphics[width=18cm, height=10cm]{images/circuit.jpg}
\end{center}
\section*{\textbf{Question Analysis:}}
\begin{itemize}
        \item From above Kmap, the output F is as follows
	\item $F=\bar{X}.\bar{Y}+Y.Z$
        \item The Truth table for above expressions is as follows
\end{itemize}
\newpage
\section*{\textbf{The Truth Table}}
\begin{center}
        \begin{table}[h]
        \input{tables/table}
\end{table}
\end{center}
\section*{\textbf{Hardware Implementation}}
The above problem is implemented and tested in hardware using Arduino UNO board. Here we used 3i/ps and 7474 IC to glow LED (output F) as per truth table and verified the expression.
\section*{Required Components \& Pin Connections}
\begin{center}
\begin{minipage}{0.45\textwidth}
\begin{table}[H]
\centering
\begin{tabular}{|c|l|}
\hline
\textbf{S.No} & \textbf{Component} \\ \hline
1 & Arduino Uno Board \\
2 & Breadboard \\
3 & 7474 IC (1) \\
4 & LEDs (1) \\
5 & Resistors: 220$\Omega$ (1) \\
6 & Jumper Wires \\
7 & USB Cable \\
\hline
\end{tabular}
\end{table}
\end{minipage}
\hspace{0.05\textwidth}
\begin{minipage}{0.45\textwidth}
\begin{table}[H]
\centering
\begin{tabular}{|c|c|}
\hline
\textbf{Component} & \textbf{Arduino Pin} \\ \hline
Input Z (5V or GND) & Digital 2 \\
Input Y (5V or GND) & Digital 3 \\
Input X (5V or GND) & Digital 4 \\
Output F (7474 D1) & Digital 8 \\
Output clk (7474 clk) & Digital 13 \\
GND & GND \\
VCC & 5V \\
\hline
\end{tabular}
\end{table}
\end{minipage}
\end{center}
\section*{Logic Description}
\begin{itemize}
        \item Let initialize inputs $X=0,Y=0,Z=0$
        \item The output expressions from KMAP and is as follows
	\item $F=\bar{X}.\bar{Y}+Y.Z$
	\item Change inputs and note Down the Experimental Truth table
\end{itemize}
\section*{Code Uploading Steps}
\begin{enumerate}
        \item Create a 7474 folder and write Makefile
        \item Write The code in main.c
        \item Run the command "make". It will compile the code and creates .hex file
        \item Copy the .hex file to ArduinoDriod folder
        \item connect the Arduino UNO to mobile with OTG cable
        \item Upload the hex file using "upload precomplied" option
        \item Observe the ouput and verify the expression
\end{enumerate}
\section*{Experimental Truth Table}
\begin{table}[H]
\centering
\begin{tabular}{|c|c|c|c|}
\hline
X & Y & Z & F (Output) \\ \hline
0 & 0 & 0 & 1 \\
0 & 0 & 1 & 1 \\
0 & 1 & 0 & 0 \\
0 & 1 & 1 & 1 \\
1 & 0 & 0 & 0 \\
1 & 0 & 1 & 0 \\
1 & 1 & 0 & 0 \\
1 & 1 & 1 & 1 \\
\hline
\end{tabular}
\end{table}
\begin{center}
        \includegraphics[width=12cm, height=12cm]{images/output.jpg}
\end{center}
\section*{Conclusion}
\begin{itemize}
        \item From Experimental Truth Table and actual truth table $F=\bar{X}.\bar{Y}+Y.Z$
        \item This matches option (B) from the original GATE question.
        \item The hardware experiment confirms the circuit's theoretical logic.
\end{itemize}

\end{document}
